\documentclass[12pt]{article}
\usepackage[T1]{fontenc}
\usepackage[utf8]{inputenc}
\usepackage{lmodern}
\usepackage{textcomp}
\usepackage{enumitem}
\usepackage{geometry}
\usepackage{graphicx}
\usepackage{amsmath, amssymb}
\geometry{margin=1in}
\begin{document}
\begin{center}
History - Undergraduate Level Assessment 2 \
Total Marks: 40 \
Answer all questions.
\end{center}

\section*{Multiple Choice (10 marks)}
\begin{enumerate}[label=\arabic*.]
\item Clovis points are NOT found in:
\begin{enumerate}[label=(\alph*)]
    \item Canada.
    \item Mexico.
    \item the continental United States.
    \item Alaska.
\end{enumerate}
\item The most prominent feature of the food resource base of post-Pleistocene Europe was its:
\begin{enumerate}[label=(\alph*)]
    \item dependence on megafauna.
    \item hunting and gathering.
    \item dependence on fur-bearing animals.
    \item diversity.
\end{enumerate}
\item During the Pleistocene era, the islands of Java, Sumatra, Bali, and Borneo formed
\_\_\_\_\_\_\_\_\_\_\_. The landmass connecting Australia, New Guinea, and Tasmania was called
\_\_\_\_\_\_\_\_\_\_\_.
\begin{enumerate}[label=(\alph*)]
    \item Sunda; Sahul
    \item Sunda; Wallacea
    \item Gondwana; Beringia
    \item Sunda; Gondwana
    \item Beringia; Sunda
\end{enumerate}
\item The undersea chasm between New Guinea/Australia and Java/Borneo is called the \_\_\_\_\_\_\_\_\_\_.
\begin{enumerate}[label=(\alph*)]
    \item Java Trench
    \item Sunda Passage
    \item Australia Trench
    \item Wallace Trench
    \item New Guinea Strait
\end{enumerate}
\item What is the name of the lithic technology seen in the Arctic and consisting of wedge-
shaped cores, micro-blades, bifacial knives, and burins?
\begin{enumerate}[label=(\alph*)]
    \item Ahrensburg Complex
    \item Solutrean Complex
    \item Aurignacian Complex
    \item Nenana Complex
    \item Gravettian Complex
\end{enumerate}
\end{enumerate}

\section*{True / False (10 marks)}
\textit{Answer True or False}
\begin{enumerate}[label=\arabic*.]
\setcounter{enumi}{5}
\item True or False: Minoan culture shows little evidence of conspicuous consumption by elites compared to other early civilizations.
\item True or False: The power of Ancestral Pueblo elites emerged significantly when they stored food surpluses for distribution in times of scarcity.
\item True or False: Sunflower was a crucial component of the Hopewell diet alongside beans and squash.
\item True or False: The tradition of making stone tools from chipped pebbles, known as Hoabinhian, originated in North America.
\item True or False: The markings on the mud-bricks of Huaca del Sol were instructions for the builders written by the engineers.
\end{enumerate}

\section*{Long Form Response (20 marks)}
\textit{Answer each question with a clear and complete explanation.}
\begin{enumerate}[label=\arabic*.]
\setcounter{enumi}{10}
\item Discuss the similarities in locomotion between Homo sapiens and Australopithecus
afarensis, and analyze how these traits may have influenced their survival and adaptation.
\item Discuss the factors that contributed to New Zealand being one of the last regions
inhabited by humans, and analyze the implications of this late settlement on its history
and culture.
\end{enumerate}

\vfill
\noindent\textit{End of Paper}
\end{document}